% Section: Evaluation protocol (drafted 2026-10-30). Every claim cites its source in a comment; every number is
% \pv{...}.
\section{Evaluation protocol}
\label{sec:protocol}

\subsection{Pre-registration and the decision record}
\label{sec:prereg}

% PREREGISTRATION.md (written before results, never edited; deviations go to DECISIONS.md: CLAUDE.md working rules);
% DECISIONS.md; docs/decision_timeline.md (R5-8); D141 (family H not run).
The project's hypotheses were pre-registered in a file written before any result and never edited afterwards.
Every experiment after that was specified in a dated entry of a public decision log, pushed to the repository before
its data existed; the commit is the timestamp, and the few amendments made after a result are labelled
(Section~\ref{sec:deviations}).
The log has \pv{record.dentries} entries.
Each pre-registration states the hypothesis, the estimator, the decision rule and what would count as failure, and
labels the analysis confirmatory or exploratory.
% Overnight working rules (decisions taken during unattended runs are logged with the alternatives considered).
Decisions taken during unattended runs were logged in the same way, with the alternatives considered.
% CLAUDE.md working rules: negative results with the same prominence; FINDINGS.md numbering.
Results are recorded as numbered findings (\pv{record.findings} so far), each citing its entry, and a negative
result is reported with the same prominence as a positive one.
% D141.
Pre-registered hypotheses that were not run are declared as such rather than dropped.
A chronological timeline of every pre-registration and amendment, with its status, accompanies the paper.

\subsection{Frozen evaluation sets}
\label{sec:frozen}

% D113 sec. 8 (R4-5: fresh seeds, manifests, FROZEN.md, the loader guard; "no model is evaluated on it ... until a
% 3-seed confirmatory run has its own entry"); D118 sec. 2 (every problem regenerated and compared before
% generation); D118 sec. 10 (interim reports: operational counts only; the analysis runs once).
The development pair sets of Section~\ref{sec:pairs} were used to build and debug every analysis.
Confirmatory claims rest on frozen sets built with fresh seeds after their analysis was specified: the R4-5 set
(A6, A6b and A6c, \pv{frozen.aSixC.problems} problems each) and the J set.
Each frozen file has a SHA-256 manifest; the evaluation loader refuses a frozen file unless the run sets an explicit
confirmatory flag, and regenerates every problem from its seed and compares it before any generation.
No model saw a frozen set before the confirmatory run that its entry registered.
While those runs were in progress, interim reports gave only operational counts (completions, truncation, parse
failures, memory backoffs), and each confirmatory analysis ran once, after all of its evaluations existed.

% CLAUDE.md component 6 (tape and replay); D72, D77 (R2-1: training-time fingerprint and the reloaded-model check);
% D78 (pinned base); D142 (resumable evals: a partial file keyed to the eval's identity).
Every completion is logged with its problem, verifier breakdown and provenance, and replaying a log regenerates each
problem from its seed and re-scores the completion exactly.
Every adapter carries a fingerprint recorded at the end of training, and an evaluation reloads it and checks the
fingerprint before generating.
Long evaluations write each finished batch to disk under a key of the model, adapter, prompts and generation settings,
so an interrupted greedy evaluation resumes without changing any completion.

\subsection{Statistics}
\label{sec:stats}

% D50/A6, D95 (the pair slope); D59 (A6b on the change in the truth); configs/pairs.yaml analysis.sigma_floor_rel.
\paragraph{Pair slopes.}
For each pair set the statistic is the within-pair OLS slope with an intercept (Section~\ref{sec:pairs}): of
$\Delta\ln\sigma_{\text{stated}}$ on $\Delta\ln\sigstar$ for A6, A6c and the J sensitivity pairs, and on
$\Delta\ln y$ for A6b and the J invariance pairs.
A stated $\sigma$ of zero enters at a relative floor of \pv{stat.floor}.
% D113 sec. 4, D118 sec. 5 (pair sets per source).
For an answer's $\sigma$, the noise-allocation slopes use pairs whose two members are both answered, and the
invariance slopes use pairs whose two members are both computation-correct, so that a wrong answer cannot pass as
invariance; values read from a propagation line use every pair.

% D104 (stage B's joint bootstrap), D118 sec. 5; Davison & Hinkley (1997).
\paragraph{Uncertainty over pairs and seeds.}
Confidence intervals come from a bootstrap over pairs with \pv{stat.B} resamples.
With several seeds, each resample draws the pairs once and recomputes every seed's statistic on them; the estimate is
the mean of the per-seed statistics~\citep{davison1997bootstrap}.
One-sided $p$-values are the share of resamples at or below zero.
% D113 R5-R8, D118 K5-K8 (expected-score differences clustered over pairs); R5-8 (never 5,390 pairs or 3 seeds as
% independent); D74, FINDINGS #22 (within-problem tests clustered over problems).
Differences in expected score (Section~\ref{sec:expected}) are averaged over seeds per problem and bootstrapped with
pairs as clusters.
Within-problem comparisons of correctness signals are clustered over problems.
Neither pairs nor seeds that share problems are ever treated as independent.

% D113 sec. 4 (C1-C3), D118 sec. 6 and D125 (P = K1 and C2 and C3); D147 sec. 8 (P_J); Schuirmann (1987) TOST;
% stage_r4.yaml analysis; pairs.yaml invariance_analysis.
\paragraph{The propagation criterion.}
A $\sigma$ source propagates if three conditions hold:
the noise-allocation slope is positive (C1: the lower end of its \pv{stat.cilevel} interval is above zero);
it is at least \pv{stat.practical}, a practical threshold (C2);
and the invariance slope is equivalent to zero, with its \pv{stat.cilevel} interval inside
$\pm\pv{stat.margin}$, a two one-sided tests criterion~\citep{schuirmann1987} (C3).
The primary hypotheses of the confirmatory runs (P on the R4-5 set, P$_J$ on the J set) apply this criterion to
D-pos-v2's answer $\sigma$, with C1 replaced by a one-sided test that is Holm-adjusted within its family.

% Holm (1979); D118 sec. 7 (K1-K8), D147 sec. 9 (J1-J6), D118 sec. 8 and D147 sec. 10 (subgroups outside the family);
% R5 central message (failure to reject is not equivalence).
\paragraph{Multiplicity.}
Each confirmatory run fixes its family of tests in advance (K1--K8 on the R4-5 set, J1--J6 on the J set), and $p$-values
are Holm-adjusted~\citep{holm1979} within it at level \pv{stat.alpha}.
Subgroup estimates (by depth, chain type or distance from the training range) are reported with intervals but are
outside the family, and no confirmatory claim rests on them.
A score difference that is not significant is reported as not rejected, never as a tie.

\subsection{Deviations}
\label{sec:deviations}

% docs/decision_timeline.md (generated Table tab:deviations); D54, D57 (the gate), D92 (X2 robustness, after the
% result), D114 (token budget; dropped rows), D125 (P amended before any confirmatory result), D131 (an eval rerun
% after a driver reset), D135 (fallback bounds chosen after seeing #35), D146 (the cross-model NO-GO accepted).
Table~\ref{tab:deviations} lists every amendment to a pre-registered plan and every deviation that touched a
registered analysis, with when it was made relative to the data it affects.
Most were made before any affected result existed.
Those made after a result are labelled as such wherever the result is reported.

% GENERATED by scripts/paper/tables.py; do not edit. Rebuild: uv run python scripts/paper/build.py
% artifact: table:deviations
% sources (path @ git blob):
%   docs/decision_timeline.md @ 5f80e454add5
\begingroup\footnotesize
\begin{longtable}{@{}l l L{7.8cm} L{3.9cm}@{}}
\caption{Amendments and deviations from the pre-registered plans.}\label{tab:deviations} \\
\toprule
entry & date & what & timing \\
\midrule
\endfirsthead
\toprule
entry & date & what & timing \\
\midrule
\endhead
\midrule
\multicolumn{4}{@{}p{15.5cm}@{}}{Every amendment and deviation in the decision timeline (dates IST, 2026), with when it was made relative to the data it affects. Source: docs/decision\_timeline.md.} \\
\endlastfoot
D47 & 10-14 10:41 & P's pre-registered sanity check failed; exploratory baselines added afterwards & amendment, \textbf{after} P's result \\
D49 & 10-14 11:54 & Two post-hoc calibration baselines; RL calibration judged against the text-feature one & amendment, before RL calibration results \\
D50 & 10-14 11:59 & Review round 1, A1--A8: two estimands; the $\sigma$$^{*}$ definition; the generator self-consistency gate (A3); statistics; nested A5 baselines and family N \ldots{} & amendment, before any result it affects \\
D54 & 10-14 12:43 & Noise capped at 5\%; 10\% becomes a heavy-tail regime & amendment, \textbf{after} the gate failed (\#9) \\
D57 & 10-14 12:56 & A3's cell-share rule replaced by a miss-count test; the next rerun is FINAL & amendment, \textbf{after} \#11, before the final rerun \\
D61 & 10-16 14:34 & A6 slope labels ("partial / full propagation") & amendment, before any A6 model result \\
D62 & 10-16 19:12 & Split-credit seeds 1 and 2; the seed-mean S1--S3 test; greedy vs sampled $\sigma$ & amendment to family S, before seed-1/2 data \\
D64 & 10-17 21:57 & The label "anti-propagation" for a CI below 0 & amendment, before any A6 model result \\
D67 & 10-18 00:06 & The bf16 merge was a bug: corrected (unmerged, 4-bit) evals become PRIMARY for families S, N, E; merged results kept, labelled superseded & amendment, before any corrected result \\
D69 & 10-18 09:48 & Post-hoc baselines re-fitted on 4-bit SFT calibration completions & amendment, before any such completion \\
D71 & 10-18 09:50 & Review round 2, R2-1--R2-7: eval hygiene (reload fingerprint), SFT revalidation, seed-level statistics (R2-4), uniform-credit seeds (R2-6), \ldots{} & amendment, before any result it affects \\
D77 & 10-19 10:31 & R2-1: below the log-prob floor, the $\sigma$-token criterion decides; the check is enforced & amendment (future fingerprints only) \\
D91 & 10-20 23:22 & X4: a direct mode, a perturbed control, the cluster bootstrap, token cost & amendment, before any X4 data \\
D92 & 10-20 23:38 & X2 robustness checks, defined before computing & analysis plan, \textbf{after} X2's result \\
D93 & 10-20 23:45 & Review round 3, R3-1--R3-10: X4 additions; X2 robustness; optimization pilots; Phase 12 scope = aleatoric propagation (R3-4); A6c (R3-5); \ldots{} & amendment, before any result it affects \\
D112 & 10-25 17:13 & Review round 4, R4-1--R4-9: D-pos-v2 vs B-v2 on mixed-level data, MC-A, the frozen set, framing (R4-8) & amendment, before any result it affects \\
D114 & 10-25 17:31 & Two deviations found while building: eval budget 1,280 tokens; 6 rows dropped & deviation, before any R4 output \\
D116 & 10-25 19:07 & D-pos-v2's eval OOM'd in its backoff; fixed and rerun at batch 16 & deviation, before any completion was scored \\
D122 & 10-26 07:52 & MC-B: two deviations found while building & deviation, before any output \\
D125 & 10-26 08:34 & D118 \S{}6 amended: the primary P is D113's full criterion (K1 $\wedge$ C2 $\wedge$ C3) & amendment, before any confirmatory result \\
D131 & 10-28 02:52 & D-pos-v2 seed 1's eval hung after a GPU driver reset; rerun on identical code & deviation (operational), before the analysis \\
D133 & 10-28 10:04 & Review round 5, R5-1--R5-10: framing; tail autopsy (R5-2); hybrid oracle (R5-3); RSS fallback (R5-4); cross-model replication (R5-5); cuts (R5-7); \ldots{} & amendment, \textbf{after} \#32 / \#33, before anything it asks for \\
D135 & 10-28 10:24 & R5-4's RSS-fallback bounds, $\sigma$/RSS $\in$ [0.333, 5], chosen on the development sets and frozen before one application to the confirmatory logs & analysis plan; grid designed \textbf{after} seeing \#35 \\
D138 & 10-28 18:46 & R2-1's fingerprint for the Llama base records at sampling batch 4 & deviation (implementation), before any pilot output \\
D139 & 10-29 00:13 & D137's go rule runs as soon as the D-pos-v2 eval exists (same code, same logs); SFT seed 1 only on GO & amendment (timing), before the go rule runs \\
\end{longtable}
\endgroup

